% ═══════════════════════════════════════════════════════════════════════════════
% CASO 1: CAPTURA INSTITUCIONAL Y MACROCRIMINALIDAD - INSTITUTO ATENEO DE COLIMA
% Estudio Documental: Genaro García Luna (2006-2012)
% ═══════════════════════════════════════════════════════════════════════════════

\nota{
  Este capítulo analiza la captura institucional del aparato de seguridad pública 
  mediante evidencia judicial (DOJ 2024), auditorías (ASF) y doctrina académica 
  (Astorga 2015; Vázquez Valencia 2019). La reconstrucción se basa exclusivamente 
  en fuentes públicas oficiales verificables.
}

\section{Ascenso técnico e inicio de la colusión (1989--2006)}

La reconstrucción biográfico-institucional demostró que la red de poder de García Luna no se originó en diciembre de 2006, sino que se consolidó de manera progresiva a lo largo de más de dos décadas dentro del aparato mexicano de inteligencia y seguridad. Ingeniero mecánico egresado de la Universidad Autónoma Metropolitana, ingresó en 1989 como agente en el Centro de Investigación y Seguridad Nacional (CISEN), donde se especializó en contraterrorismo y combate al secuestro y articuló vínculos estrechos con mandos políticos, policiales y militares \citep{castillog2023,expansion2019}.

En 1999 se incorporó a la recién creada Policía Federal Preventiva y en diciembre de 2000 asumió la dirección de la Policía Judicial Federal. Durante el gobierno de Vicente Fox (2001--2006), diseñó y encabezó la Agencia Federal de Investigación (AFI), proyectando una reputación internacional de modernizador policial avalada por reconocimientos de INTERPOL y agencias estadounidenses, aun frente a controversias públicas graves como el montaje televisivo del caso Florence Cassez en diciembre de 2005 \citep{astorga2007,expansion2019}. No obstante, de acuerdo con las pruebas acreditadas posteriormente ante la justicia federal de Estados Unidos, fue desde el año 2001 cuando comenzaron sus acuerdos de soborno y colusión con el Cártel de Sinaloa \citep{doj2024garcialuna,pastrana2024}. Este hallazgo coincide con lo señalado por Flores Pérez (2009), quien documentó cómo la alternancia política sin controles efectivos sobre los cuerpos de policía permitió la infiltración de redes criminales en instituciones de seguridad.

\section{Concentración sin contrapesos y doble vertiente de corrupción (2006--2012)}

Al iniciar su mandato, Felipe Calderón nombró a García Luna titular de la Secretaría de Seguridad Pública (SSP), otorgándole una concentración inédita de presupuesto, inteligencia y mando operativo en el marco de la estrategia contra el narcotráfico \citep{astorga2015}. Bajo su gestión, el estado de fuerza de la Policía Federal se expandió de alrededor de 5,000 a más de 35,000 efectivos al cierre del sexenio \citep{pastrana2024}. Dicha concentración discrecional ---en ausencia de contrapesos horizontales efectivos \citep{lopezayllon2009,zaid2019}--- se asoció con la operación simultánea de dos vertientes de gran corrupción:

\subsection{Protección criminal (cohecho y abuso de funciones)}

En el expediente 19-CR-576 (S-1) (BMC) sustanciado en la Corte Federal del Distrito Este de Nueva York, la fiscalía acreditó que el exsecretario recibió decenas de millones de dólares en sobornos en efectivo del Cártel de Sinaloa. A cambio, facilitó el tráfico de más de un millón de kilogramos de cocaína hacia Estados Unidos, filtró información de inteligencia sobre operativos en curso y dirigió acciones policiales contra organizaciones competidoras \citep{doj2024garcialuna,insightcrime2023}. Al dictar sentencia, el juez federal Brian Cogan resumió la magnitud institucional del caso al señalar que el acusado llevó una «doble vida» equiparable a la de los capos que combatía públicamente \citep{pastrana2024}. En el plano operativo interno, su colaborador cercano y exdirector de Seguridad Regional de la Policía Federal, Luis Cárdenas Palomino, fue detenido en México en 2021 por el delito de tortura y figuró en las indagatorias estadounidenses.

\subsection{Red contractual y lavado de activos (peculado y desvío de recursos)}

De forma paralela a los sobornos del narcotráfico, a partir de 2009 ---con continuidad transexenal hasta 2018 en contratos con el CISEN, la Policía Federal, el Órgano Administrativo Desconcentrado Prevención y Readaptación Social (OADPRS) y la entonces Procuraduría General de Justicia de la Ciudad de México--- García Luna articuló junto con la familia empresarial Weinberg (Mauricio Samuel y Jonathan Alexis Weinberg) una estructura corporativa offshore registrada en Panamá e Israel (Nunvav Inc., Nunvav Technologies, Nice Systems Ltd. y la consultora GLAC) \citep{mcci2023,castilloj2026}. La Unidad de Inteligencia Financiera (UIF) documentó 30 contratos públicos por más de 727 millones de dólares y 528 millones de pesos para equipamiento tecnológico y mantenimiento penitenciario, cuyos fondos fueron dispersados hacia cuentas en Barbados, Israel y México y hacia empresas factureras identificadas por el SAT \citep{uif2024,mcci2023}.

\section{Judicialización internacional, sentencias y posturas de descargo (2019--2026)}

Aunque diversas denuncias públicas y señalamientos internos se habían formulado años atrás sin traducirse en procesos penales firmes dentro de México, la ruptura judicial se produjo en el extranjero. Tras ser arrestado por agentes de la DEA en Dallas, Texas, en diciembre de 2019, Genaro García Luna fue declarado culpable de cinco cargos federales vinculados al narcotráfico y falsedad de declaraciones en febrero de 2023, y condenado el 16 de octubre de 2024 a más de 38 años de prisión, sentencia que comenzó a purgar en la penitenciaría federal de máxima seguridad ADX Florence \citep{doj2024garcialuna}.

En la vía civil patrimonial, la demanda promovida por el Estado mexicano mediante la UIF ante cortes de Florida obtuvo resoluciones condenatorias de alta cuantía: el 22 de mayo de 2025, un tribunal de Miami ordenó a García Luna y a su esposa, Linda Cristina Pereyra, resarcir aproximadamente 2,500 millones de dólares al Estado mexicano (1,740 millones correspondientes a Pereyra) \citep{uif2025,guillen2025}. Posteriormente, el 17 de noviembre de 2026, la Secretaría de Hacienda y Crédito Público notificó la obtención de una nueva sentencia civil favorable a México por más de 578 millones de dólares dentro del mismo entramado empresarial \citep{uif2026}.

En cumplimiento del criterio de confirmabilidad y balance probatorio, se registró la postura de descargo del expresidente Felipe Calderón, quien tras la sentencia de octubre de 2024 rechazó toda responsabilidad personal y declaró que durante su gestión nunca dispuso de «evidencia verificable que lo involucrara con actividades ilícitas» ni recibió alertas formales de agencias de inteligencia nacionales o estadounidenses, defendiendo la legalidad de su estrategia de seguridad \citep{infobae2024}. Sin embargo, el contraste entre las sentencias judiciales estadounidenses, la continuidad contractual hasta 2018 y la ausencia de sanción penal interna oportuna evidenció debilidades estructurales profundas en los mecanismos de supervisión horizontal y bilateral \citep{insightcrime2023,merino2024,vazquezvalencia2019}.